\section*{Your turn}

Everything here needs three processes and Docker. Check out tag \code{ch10} of
the companion repository, run \code{npm ci} once, and from the repository root
start the database, both subgraphs and the router:

\begin{minted}[fontsize=\footnotesize,breaklines]{text}
docker compose up -d mosaic-db
dotnet run --project src/Mosaic.Catalog     # 5101
dotnet run --project src/Mosaic.Api         # 5100
docker compose up -d mosaic-router          # 3002
\end{minted}

The two \code{dotnet run} commands want a terminal each. If port 3002 is
already taken, an earlier run of \code{docker compose --profile wire up} left
chapter~\ref{ch:how-a-federated-query-actually-runs}'s router behind;
\code{docker compose --profile wire down} clears it. The commands below are
written for PowerShell, and the only thing that changes under \code{sh} is the
line continuation.

\begin{enumerate}
  \item \textbf{Answer the question the book could not.} Send
  section~\ref{sec:ch10-the-query-that-stopped-working}'s storefront query to
  \code{http://localhost:3002/graphql}, then send the same document to 5100 and
  to 5101. You get one answer and two 400s, and the two failures name different
  fields. Keep the three responses side by side for a moment. They are the
  argument for the whole of Part II on one screen.

  \item \textbf{Watch a plan get smaller.} Ask for the plan of the storefront
  query with \code{X-WG-Include-Query-Plan} and \code{X-WG-Skip-Loader}, then
  remove \code{price} from the selection, then \code{availableQuantity}, then
  \code{averageRating}, asking for the plan each time. The
  \code{dependencies} array shrinks by one entry each time and the number of
  fetches does not change. Now remove \code{reviews} as well: the second fetch
  disappears entirely. That is section~\ref{sec:ch10-the-plan-on-a-real-graph}'s
  point about what the array's length is a count of, felt rather than read.

  \item \textbf{Predict the plan for two root fields.} Before running it, write
  down how many fetches you expect from

  \begin{minted}[fontsize=\footnotesize,breaklines]{text}
{ browseProducts(first: 2) { nodes { title } } customerById(id: "...") { displayName } }
  \end{minted}

  with a real customer identifier, which you can get from a review's author.
  Then run it. If your prediction was two fetches under a \code{Sequence}, look
  again at what node they are under and why
  chapter~\ref{ch:how-a-federated-query-actually-runs} says the router
  serialises only as much as the data forces it to.

  \item \textbf{Make the literals matter, and find that they do not.} Send the
  storefront query at \code{first: 3}, then at \code{first: 7}, capturing both
  plans. Compare them as strings rather than by eye. Then send the same
  document with proper variable declarations and compare again. Two of the
  three comparisons come out equal and one does not, and the one that does not
  is the one a generated client sends.

  \item \textbf{Break the routing on purpose.} Copy the router's configuration
  to a second file beside it, and add this line to the copy:

  \begin{minted}[fontsize=\footnotesize]{text}
localhost_fallback_inside_docker: false
  \end{minted}

  Start a second router on it. It needs a port of its own, because
  \code{mosaic-router} already holds 3002:

  \begin{minted}[fontsize=\footnotesize,breaklines]{text}
docker run --rm -d --name nofallback -p 3003:3002 `
  --add-host host.docker.internal:host-gateway `
  -v ${PWD}/federation/supergraph.json:/etc/mosaic/supergraph.json:ro `
  -v ${PWD}/router/nofallback.yaml:/etc/mosaic-router/config.yaml:ro `
  -e CONFIG_PATH=/etc/mosaic-router/config.yaml `
  ghcr.io/wundergraph/cosmo/router:0.337.1
  \end{minted}

  Send the storefront query to port 3003. Note the status code before you read
  the body, then decide what a load balancer in front of this would have
  concluded. \code{docker rm -f nofallback} puts things back.

  \item \textbf{Set a value in two places.} Change \code{log\_level} to
  \code{debug} in \code{router/config.yaml}, then recreate the container:

  \begin{minted}[fontsize=\footnotesize,breaklines]{text}
docker compose up -d --force-recreate mosaic-router
docker compose logs mosaic-router
  \end{minted}

  Confirm the debug lines appear. Now put the file back to \code{info} and set
  the environment variable instead, by adding this under \code{mosaic-router}
  in \code{docker-compose.yml}:

  \begin{minted}[fontsize=\footnotesize]{yaml}
      LOG_LEVEL: "debug"
  \end{minted}

  Recreate it again and count the debug lines. Then read the second line of the
  startup log, which told you this would happen. Put both files back with
  \code{git checkout}.

  \item \textbf{Reload the graph under a running router.} With the router up,
  edit a field description in \code{src/Mosaic.Catalog}, re-export that
  service's schema, and recompose over the committed config. Watch the router's
  log for \code{Config file changed}, then introspect to confirm the
  description changed. Nothing restarted. Now run the verification script: it
  fails, because the committed config no longer matches the one your edit
  produced. Put both back with \code{git checkout}.

  \item \textbf{Deploy a graph you were told not to.} Run

  \begin{minted}[fontsize=\footnotesize,breaklines]{text}
node scripts/router-cases.mjs --print resolvability-off
  \end{minted}

  and read the five lines of evidence it prints. Then do it by hand:

  \begin{minted}[fontsize=\footnotesize,breaklines]{text}
# add ", resolvable: false" to the @key on Product in schema/mosaic.graphql
npx wgc router compose --disable-resolvability-validation `
    -i federation/mosaic.yaml -o federation/scratch.json
  \end{minted}

  Then start a router on that file, on a port of its own again, with the
  \code{docker run} of exercise 5 and \code{scratch.json} in place of
  \code{supergraph.json}. Try to find out what is wrong with your graph using
  only what the router tells a client. Give yourself five minutes. Then look at
  \code{docker logs}. Afterwards, restore the schema with
  \code{git checkout schema/mosaic.graphql}.

  \item \textbf{Ask the router for something only a subgraph has.} Send
  \code{\{ \_service \{ sdl \} \}} to 3002 and to 5101. Then send
  \code{\{ nope1 nope2 nope3 \}} to both. The status codes differ, the wording
  differs, and one of them reports three errors where the other reports one.
  Decide which behaviour you would rather your clients see, and note that you
  do not get to choose.
\end{enumerate}

Exercise 8 is the one worth doing properly, with the five minutes actually
spent. Reading that a diagnostic is useless is not the same experience as
needing it and finding it is not there, and the second one is what makes you
put composition in a pipeline the following Monday.
Chapter~\ref{ch:entity-resolution-done-right} takes the second hop apart: what
a reference resolver does when the router batches twenty-five representations
into it, and what \code{@requires} costs when a field needs something the other
service owns.
